\documentclass[a4paper]{article}
\usepackage[a4paper, total={6in, 9in}]{geometry}
\usepackage[english]{babel}
\usepackage{enumitem}
\usepackage{amsmath}
\usepackage{amsthm}
\usepackage{amssymb}
\newtheorem{theorem}{Theorem}
\newtheorem{corollary}{Corollary}[theorem]
\newtheorem{lemma}[theorem]{Lemma}
\newtheorem{definition}{Definition}[section]
\usepackage{graphicx}
\usepackage{listings}
\usepackage{hyperref}
\usepackage{algorithm}
\usepackage[noend]{algpseudocode}
\usepackage[svgnames]{xcolor}

\newcommand{\instruction}[1]{\textcolor{blue}{#1}}

\title{Practical 1: Festival Road Closings}								
\author{Lan Stare}

\makeatletter
\let\thetitle\@title
\let\theauthor\@author
\let\thedate\@date
\makeatother


\begin{document}

%%% Title %%%
\begin{center}
\rule{\linewidth}{0.2 mm} \\[0.4 cm]
{ \huge \bfseries \thetitle}\\
\rule{\linewidth}{0.2 mm} \\[0.5 cm]
\textsc{\Large Algorithms and Data Structures}\\[0.2 cm]				
\textsc{\large Practical 1}\\[0.5 cm]				

\begin{minipage}{0.4\textwidth}
    \begin{flushleft} \large
        \emph{Group number: 42}\\
        \theauthor\
        \end{flushleft}
        \end{minipage}~\
        \begin{minipage}{0.4\textwidth}
        \begin{flushright} \large~\\~\\	s1169977
    \end{flushright}
\end{minipage}\\[1 cm]
\end{center}


\section{Explanation of the Algorithm}
\subsection{Initialization}
\subsubsection*{Reading in the data}
We read in the data by first saving the entire input as a list of integers using the sys library. This allows us to easily split the flattened row into integers regardless of the format of the \textsl{.in} files, since some files only contain the number of vertices ($|V|$) and edges ($|E|$) in the first line (and every subsequent line contains exactly one edge) while others contain all the data in a single line. After reading in the data, we extract $|V|$ and $|E|$ from the list of integers we had just created. Finally, we build a list of edge tuples by iterating through the list of integers and grouping every three integers into a tuple representing an edge ($v1$, $v2$, $type$).

\subsubsection*{Initializing data structures}
The main data structure we use to store our graph is a list of tuples, where each tuple represents an edge in the graph. Each tuple contains three elements: the starting vertex, the ending vertes, and the type of edge we have. Type 0 repredents a pedestrian-only road, type 1 represents a bus-only road and type 2 represents a road that can be used by both pedestrians and busses (later referred to as a \underline{combined road}). We use this data structure because it allows us to easily iterate through all the edges and we also don't need to keep track of adjacent vertices, which would be possible with an adjacency list or matrix but would be more expensive. Duplicate edges of different types are also possible in the input, so dictionaries are not suitable for this problem.


\subsection{Algorithm Logic}

\subsubsection*{Algorithm: Initialization}
The core idea is using \underline{Kruskal's algorithm}. Since we haven't separated the graph into a pedestrian and bus graph prior to running the algorithm, we have to take care of both (or all three) types of edges while iterating.

We first build two union-find data structures, one for the pedestrian graph and one for the bus graph. This allows us to keep track of connected components for both types. We then need to build a find function that finds the root of a vertex in the union-find data structure and a union function that connects two vertices in the union-find if they are not already in the same component (when their roots are different).

Before iterating through the edges, we need to sort them by their type. We want to prioritize the combined edges (type 2), so we assign them weight 1 while assigning weight 2 to the pedestrian-only and the bus-only edges (types 0 and 1). 

Finally, we initialize a variable \textsl{kept\_edges} to 0, which will keep track of the number of edges we keep in the final graph.

\subsubsection*{Algorithm: Main loop}
Up intil now, we have only implemented the initiation of the Kruskal's algorithm. The main loop is where we execute our modifications to the original algorithm.

First we add a switch which is set to False in the beginning and will keep track of whether we have added an edge in the current iteration or not. We then iterate through the sorted list of edges. If the road is of type 1 or 2 (bus-only or combined), we check if the two vertices of the edge are in different components in the bus union-find data structure. If they are, we connect them and set the switch to True. We do the same if the road is of type 0 or 2 (pedestrian-only or combined). If the switch is set to True after checking both if statements, we increment the \textsl{kept\_edges} by 1.

\subsubsection*{Arguments for used structures and algorithms}
We use Kruskal's algorithm because we want to find the minimum amount of edges that connect all vertices in both graphs. The graph is undirected so we could have used Prim's algorithm as well with binary heaps (to get the same time complexity) but Kruskal's algorithm is easier to implement and modify for our needs. 

Instead of only checking whether adding an edge would create a cycle in one graph, we have to check both graphs. To do this efficiently, we separated the union-find data structures for both graphs, which allows us to keep track of connected components for both graphs at the same time.

Later we also added an extra stopping criterion to the main loop that checks whether both graphs are fully connected, further explained in the next section.

\subsection{Returning the Answer}
The original algorithm terminates when we have iterated through all the edges. However, we added a stopping criterion that checks whether both the bus and pedestrian subgraphs are fully connected. We can check this by keeping track of the number of connected components in both union-find data structures. If both graphs have only one connected component left, we can stop iterating through the edges since union operations will no longer be possible. The reason for this is that once two components are connected, they will never be disconnected again, meaning that once we reach one connected component, it will stay that way (otherwise we would form a cycle).

In the end, the function has to subtract the number of kept edges from the total number of edges to get the number of removed edges. The function then returns this number as the final output.

\subsection{Optimizations}
The main optimization we made to the algorithm was checking whether the subgraphs were fully connected during iterating as opposed to checking it at the end. This both allowed us to stop iterating earlier (although not changing the time complexity in theory) and also returning -1 if a spanning tree for pedestrians or for busses didn't exist. A more expensive alternative would be counting the number of roads used in the final subgraph, while regarding the combined roads as 2 ordinary edges. If the final number of edges was not $2|V| - 2$, the algorithm would return -1.

The optimization mentioned is done by first initiating two variables \textsl{bus\_components} and \textsl{pedestrian\_components} to $|V|$ (the number of vertices). In the main loop we then decrement the corresponding variable by 1 if we connect two components. With this technique we can check in constant time whether both graphs are fully connected by checking if both variables are equal to 1. If they are not, we continue iterating through all the edges. If they still are not in the end, we return -1 since it is impossible to connect all vertices in both graphs.

Another optimization was counting the edges as they are added to the subtree as opposed to building and saving the subtree and then traversing it to check how many edges it had. This reduced the time complexity of the operation from at least $O(|E|)$ to $O(1)$. This, however, doesn't play a big role in the overall complexity of the algorithm.

\section{Correctness Analysis}
\begin{theorem}
    Double Kruskal's algorithm finds the maximum amount of roads that can be closed between the stages for both busses and pedestrians.
\end{theorem}

\begin{proof}
    Kruskal's algorithm is an algorithm that efficiently finds the minimum spanning subtree of a connected graph. If our problem was finding the maximum amount of roads that can be closed for pedestrians-only (or symmetrically busses-only), the basic Kruskal's algorithm would accurately find a spanning subtree which corresponds to minimum amount of roads that can be used, since spanning trees still have all the vertices of the original graph.

    The \textsl{double} Kruskal's algorithm in this context means that we are constructing 2 spanning trees within the same graph. This means that the resulting subgraph is generally not a spanning tree, but both of them (by itself) are.

    By weighting the combined roads as 1 and all the other roads (pedestrian or bus-only) as 2, we make sure that Kruskal's algorithm prioritizes the combined roads. In the shared graph we get in the end, the combined roads count as double since they are used in both the pedestrian and the bus spanning tree. Since Kruskal's algorithm finds the minimum spanning tree, the combined roads have to weigh half of the weight of the other roads, to be prioritized.

    We only traverse each road once in the double Kruskal's algorithm while checking whether each new road forms a cycle in either subgraphs that we had built until then. This removes the possibility of any combined road ending up in only one of the subgraphs. It also avoids adding the pedestrian and bus roads before all the combined roads used had been added.

    The stopping criterion we used in this algorithm was when both spanning trees by itself get $|V| - 1$ edges, the algorithm terminates. This is valid for Kruskal's algorithm, because a spanning subtree contains all vertices of the original graph ($|V|$) and is a tree, which means that it has to contain exactly $|V| - 1$ edges.

    By counting the roads when they are added, we efficiently calculate the number of roads used in the final subgraph. To get the maximum amount of roads that can be removed, we have to subtract the roads used from all the roads in the original graph.


\end{proof}


\section{Complexity Analysis}
\subsubsection*{Reading in the data}

\begin{theorem}
    Complexity of reading in the data is $O(|E|)$ where $|E|$ is the number of edges in the graph.
\end{theorem}
\begin{proof}
    The function \textsl{sys.stdin.read()} reads through all the characters in the input once, which takes $O(c)$ time where $c$ is the number of characters in the input. The \textsl{.split()} function splits the string on whitespace chracters, which is also done in $O(c)$ time, since it has to check each character once. The \textsl{map(int, $\ldots$)} then converts each of the ($3|E| + 2$) substrings into integers in $O(|E|)$ time. For building the tuples of edges, we iterate through the list of integers once more, making slices of 3 integers at a time (one for every edge), which takes $O(|E|)$ time. Since the file contains $3|E| + 2$ integers and $c$ is proportional to the number of integers, we can conclude that reading the data takes $O(|E|)$ time in total.
\end{proof}

\subsubsection*{Main algorithm}
\begin{theorem}
    Complexity of the main algorithm is $O(|E| \log |E|)$.
\end{theorem}
\begin{proof}
    \textbf{Initialization:} Initializing the union-find data structures (\textsl{parent\_bus = list(range(n))} and \textsl{parent\_ped = list(range(n))}) takes $O(|V|)$ time each, since creating a list of size $n$ parses through all $n$ elements once. Both initializations of the number of components clearly happen in constant time. The variable \textsl{kept\_edges} is also initialized in constant time.

    Sorting the edges in python runs in $O(|E| \log |E|)$ time.

    \textbf{Main loop:} The main loop iterates through all edges once, which takes $O(|E|)$ time. Inside the loop, we have at most one call to the union function for each edge. Furthermore, the union function consists of two calls to the find function and some constant time operations. Therefore, we only need to analyze the time complexity of the union-find data structure.

    Since the find and the union functions are implemented exactly like they were at one of our lectures, we can assume that they both run in $O(\log|E|)$ time complexity.
    
    We can see that the main comtributors towards complexity are the sorting and the main loop, making the total complexity of the algorithm: $O(|E| \log |E|)$.
\end{proof}


\section{Reflection}

\subsubsection*{Building the algorithm}
I first wrote down everything I could have used from what we had discussed in the lectures or used in the weekly assignments. I came to the conclusion of using Kruskal's algorithm pretty fast but didn't immediately know what would be the best way to implement it. I struggled quite long with building the right algorithm that would produce the right output. My first idea was to separate the graph into two subgraphs first and first iterate through the bus graph remembering the combined edges used and then running Kruskal's algorithm again on the pedestrian graph while marking the used edges in the previous iteration as \textsl{required edges} and adding them to the connected edges immediately. After getting the wrong output I realized that the idea wasn't going to work since the algorithm counted the same edges twice and was so computationally expensive that it had trouble computing the output for even a single file. I scratched the idea and thought about how to avoid running Kruskal's algorithm twice so the idea of simultanious algorithm without prior graph separation came to mind.

\subsubsection*{Speeding up the algorithm}
The last adjustment to the algorithm was speeding it up by computing the number of components which was useful to check whether both busses and pedestrians were able to reach all stages while iterating through the edges. I had previously computed it in the end by checking the number of edges in the constructed bus and pedestrian graphs. There surely exist further optimizations for my approach, for example a faster sorting algorithm or a library for python or another language that supports the (double) Kruskal's algorithm better.

\end{document}